% Options for packages loaded elsewhere
\PassOptionsToPackage{unicode}{hyperref}
\PassOptionsToPackage{hyphens}{url}
\documentclass[
  ignorenonframetext,
]{beamer}
\newif\ifbibliography
\usepackage{pgfpages}
\setbeamertemplate{caption}[numbered]
\setbeamertemplate{caption label separator}{: }
\setbeamercolor{caption name}{fg=normal text.fg}
\beamertemplatenavigationsymbolsempty
% remove section numbering
\setbeamertemplate{part page}{
  \centering
  \begin{beamercolorbox}[sep=16pt,center]{part title}
    \usebeamerfont{part title}\insertpart\par
  \end{beamercolorbox}
}
\setbeamertemplate{section page}{
  \centering
  \begin{beamercolorbox}[sep=12pt,center]{section title}
    \usebeamerfont{section title}\insertsection\par
  \end{beamercolorbox}
}
\setbeamertemplate{subsection page}{
  \centering
  \begin{beamercolorbox}[sep=8pt,center]{subsection title}
    \usebeamerfont{subsection title}\insertsubsection\par
  \end{beamercolorbox}
}
% Prevent slide breaks in the middle of a paragraph
\widowpenalties 1 10000
\raggedbottom
\AtBeginPart{
  \frame{\partpage}
}
\AtBeginSection{
  \ifbibliography
  \else
    \frame{\sectionpage}
  \fi
}
\AtBeginSubsection{
  \frame{\subsectionpage}
}
\usepackage{iftex}
\ifPDFTeX
  \usepackage[T1]{fontenc}
  \usepackage[utf8]{inputenc}
  \usepackage{textcomp} % provide euro and other symbols
\else % if luatex or xetex
  \usepackage{unicode-math} % this also loads fontspec
  \defaultfontfeatures{Scale=MatchLowercase}
  \defaultfontfeatures[\rmfamily]{Ligatures=TeX,Scale=1}
\fi
\usepackage{lmodern}
\ifPDFTeX\else
  % xetex/luatex font selection
\fi
% Use upquote if available, for straight quotes in verbatim environments
\IfFileExists{upquote.sty}{\usepackage{upquote}}{}
\IfFileExists{microtype.sty}{% use microtype if available
  \usepackage[]{microtype}
  \UseMicrotypeSet[protrusion]{basicmath} % disable protrusion for tt fonts
}{}
\makeatletter
\@ifundefined{KOMAClassName}{% if non-KOMA class
  \IfFileExists{parskip.sty}{%
    \usepackage{parskip}
  }{% else
    \setlength{\parindent}{0pt}
    \setlength{\parskip}{6pt plus 2pt minus 1pt}}
}{% if KOMA class
  \KOMAoptions{parskip=half}}
\makeatother
\setlength{\emergencystretch}{3em} % prevent overfull lines
\providecommand{\tightlist}{%
  \setlength{\itemsep}{0pt}\setlength{\parskip}{0pt}}
\usepackage{bookmark}
\IfFileExists{xurl.sty}{\usepackage{xurl}}{} % add URL line breaks if available
\urlstyle{same}
\hypersetup{
  hidelinks,
  pdfcreator={LaTeX via pandoc}}

\author{\texorpdfstring{}{}}
\date{}

\begin{document}

\section{\texorpdfstring{Conclusion: Evidence Landscape, Methodological
Limitations, and Research Agenda
\label{sec:conclusion}}{Conclusion: Evidence Landscape, Methodological Limitations, and Research Agenda }}\label{conclusion-evidence-landscape-methodological-limitations-and-research-agenda}

\begin{frame}{Summary}
\protect\phantomsection\label{summary}
This work demonstrates a first-generation prototype infrastructure for
computational meta-analysis of a rapidly growing scientific field. By
combining multi-source retrieval (\(N = 849\) papers from three
databases), LLM-based assertion extraction encoded as nanopublications,
and citation-weighted hypothesis scoring, we produce a queryable,
RDF-compatible knowledge graph that tracks the evolving evidence for
eight core Active Inference claims. The system demonstrates the
feasibility of automated living reviews, while clearly delineating the
boundaries of current model capabilities.
\end{frame}

\begin{frame}{Constraints and Methodological Scope}
\protect\phantomsection\label{constraints-and-methodological-scope}
Several conscious design constraints scope these findings.

\begin{block}{Keyword Classifier Resolution}
\protect\phantomsection\label{keyword-classifier-resolution}
The keyword-based classifier operates over 65+ mathematical indicators
distributed across 8 domain categories, using a deterministic priority
system that routes papers to specific application domains (C1--C5)
before testing tools (B), formal theory (A1), and the qualitative
philosophy catch-all (A2). Word-boundary-aware matching reduces
partial-match false positives, but keyword-based methods cannot capture
semantic nuance: papers using novel terminology or discussing
cross-domain topics without standard vocabulary risk misclassification.
Residual A2 concentration should be interpreted as a ceiling on broad
theoretical generality rather than a literal measure of philosophical
focus. An embedding-based classifier trained on a labeled subset would
provide a quantitative upper bound on the fraction of A2 papers that
merit redistribution.
\end{block}

\begin{block}{Citation Network Coverage Gaps}
\protect\phantomsection\label{citation-network-coverage-gaps}
The 1,678 intra-corpus edges spanning 567 connected components provide a
topological skeleton, but three systematic gaps inflate the component
count: (1) cross-source identifier mismatches (DOI vs.~OpenAlex
vs.~arXiv ID), (2) papers whose references are not indexed by any source
API, and (3) open-access preprints whose DOIs differ from their
published versions. Exhaustive DOI-level cross-matching with fuzzy title
matching would condense the graph further.
\end{block}

\begin{block}{Corpus Biases, Citation Dynamics, and Linguistic Framing}
\protect\phantomsection\label{corpus-biases-citation-dynamics-and-linguistic-framing}
Citation counts are subject to Matthew effects and cumulative field-size
biases. Partial-year indexing for the most recent calendar year
undercounts recent publications. The measured 16.99\% CAGR reflects the
dilutive effect of the long temporal span (2005--2026); the growth phase
from 2010 onward follows a steeper trajectory. Additionally, the
retrieved corpus itself suffers from selection biases inherent to
queried databases, including English-language dominance and the
structural over-indexing of preprints relative to peer-reviewed final
versions. Finally, the predominantly positive hypothesis scores across
the board are inflated by two systematic effects: (1)
\textbf{publication bias}, which causes academic journals to
preferentially select positive and confirmatory findings
\citep{sterling1959publication}, and (2) \textbf{linguistic asymmetry}
in scientific writing, where declarative claims are phrased
affirmatively far more often than negatively. These effects jointly
suppress contradicting assertions in the extracted evidence base.
Relative rankings and temporal trajectories are more reliable than
absolute score magnitudes.
\end{block}

\begin{block}{LLM Extraction Fidelity, Domain Drift, and Robustness}
\protect\phantomsection\label{llm-extraction-fidelity-domain-drift-and-robustness}
Zero-shot LLM extraction introduces distinct systematic biases:
over-extraction (the model hallucinating certainty for claims the paper
merely mentions in passing) and direction inversion (misclassifying
opposing evidence as supporting). Recent benchmarking confirms that
state-of-the-art systems often fall short of production-level precision
on tasks requiring exhaustive retrieval and aggregation of directional
claims from long documents \citep{liang2024survey}. Furthermore, because
our corpus extends to 2026, LLM extraction is vulnerable to \emph{domain
drift}---the base models may lack parametric knowledge of the most
recent theoretical developments. As an alternative, fine-tuned models
specifically trained on FEP/AIF abstracts could yield higher precision
than our zero-shot approach, though at a steeper computational setup
cost.

To formally bound these extraction errors and ensure robustness, we
instituted the 10\% manual-annotation ground-truth baseline evaluated
via Cohen's \(\kappa\) (detailed in \S\ref{sec:extraction_pipeline}).
Evaluating the automated scoring pipeline against this baseline ensures
it meets minimum inter-rater reliability thresholds (\(\kappa > 0.70\))
before aggregating the data. The explicit ``irrelevant'' filtering
predicate further mitigates over-extraction, converting what would be a
vulnerability of automated reviews into a calibrated evidential ledger.
\end{block}
\end{frame}

\begin{frame}[fragile]{Future Directions: Beyond Tally-Based Evidence
Aggregation}
\protect\phantomsection\label{future-directions-beyond-tally-based-evidence-aggregation}
The current scoring formula
(\hyperref[sec:methods_kg]{described in the methodology}) aggregates
LLM-extracted assertions through a simple citation-weighted tally. While
this approach provides a transparent and reproducible baseline, it
leaves substantial room for methodological sophistication. We identify
six directions, ordered by expected impact, with the first three
specifically addressing the limitations of tally-based evidence
synthesis.

\begin{block}{Hierarchical Bayesian Hypothesis Scoring}
\protect\phantomsection\label{hierarchical-bayesian-hypothesis-scoring}
The most direct extension replaces the additive tally with a
\textbf{hierarchical Bayesian model} that treats each hypothesis score
as a latent variable inferred from noisy assertion observations. Under
this formulation, each assertion \(a_i\) contributes a likelihood term
\(P(a_i | \theta_H, \sigma)\) parameterized by the hypothesis-level
evidence strength \(\theta_H\) and an observation noise term \(\sigma\)
capturing LLM extraction uncertainty. A hierarchical prior
\(\theta_H \sim \mathcal{N}(\mu_{\text{field}}, \tau^2)\) pools
information across hypotheses, enabling principled shrinkage for
hypotheses with sparse evidence (e.g., H6 Clinical Utility, with only 35
assertions). This framework produces posterior credible intervals rather
than point estimates, providing uncertainty quantification that the
current tally-based scores lack. Temporal dynamics can be modeled
through time-varying parameters \(\theta_H(t)\) using state-space
formulations that re-weight older evidence rather than treating all
cumulative assertions equally.
\end{block}

\begin{block}{Causal Evidence Graphs}
\protect\phantomsection\label{causal-evidence-graphs}
A second-generation knowledge graph would encode not only
assertion-level relationships (paper → supports → hypothesis) but also
\textbf{causal dependencies among hypotheses} themselves. For example,
evidence for predictive coding (H4) often implicitly supports FEP
universality (H1), yet the tally-based approach treats them as
independent. A causal evidence graph---structured as a directed acyclic
graph (DAG) over hypotheses with edge weights learned from co-assertion
patterns---would enable cross-hypothesis evidence propagation using
belief propagation or variational message passing. This is particularly
relevant for the Active Inference literature, where hypotheses are
theoretically nested: FEP universality (H1) logically entails predictive
coding (H4), and Markov blanket realism (H3) is a prerequisite for
certain formulations of H1. Encoding these dependencies would prevent
the double-counting of evidence from papers that support multiple
related hypotheses and enable identification of which specific claims
drive support for downstream hypotheses. The resulting causal structure
itself would be a scientific contribution---a formal map of evidential
dependencies within the field's theoretical architecture.
\end{block}

\begin{block}{Evidential Diversity and Source Weighting}
\protect\phantomsection\label{evidential-diversity-and-source-weighting}
The current formula weights assertions by
\(\log(1 + \text{citations}) \cdot \text{confidence}\), treating all
assertion sources symmetrically. A more nuanced approach would introduce
an \textbf{evidential diversity index} that downweights correlated
evidence from papers sharing authors, institutions, or methodological
approaches. Concretely, assertions could be weighted by the inverse of
their similarity to previously counted assertions, measured via cosine
similarity of paper embeddings. This would address the observation that
H1 (FEP universality) accumulates a large neutral tally partly because
many A2 (philosophy) papers invoke the FEP without independently testing
it---a form of evidential redundancy that inflates the evidence base
without adding independent information. Additionally, assertions could
be stratified by evidence type (empirical, theoretical, review) with
configurable type-specific weights, enabling users to compute evidence
scores that privilege experimental results over theoretical commentary.
\end{block}

\begin{block}{Additional Directions}
\protect\phantomsection\label{additional-directions}
\begin{enumerate}
\item
  \textbf{Confidence calibration.} A pilot study comparing LLM-generated
  assertions with domain expert assessments would establish
  inter-annotator agreement (\(\kappa\)) and identify systematic biases.
  This is the prerequisite for all downstream improvements.
\item
  \textbf{Agentic LLM Extractors.} Drawing on recent work connecting
  active inference to artificial reasoning \citep{friston2025active} and
  proposing AIF as a reward-free alternative for LLM-based agents
  \citep{wen2025missing}, replacing static prompt templates with
  goal-directed, actor-critic LLM architectures could significantly
  improve confidence calibration. The Renormalization Generative Model
  (RGM) architecture introduced by Friston et
  al.~\citep{friston2025active}---which achieves 99.8\% MNIST accuracy
  with 90\% less training data through scale-free hierarchical
  generative modeling---demonstrates that principled AIF architectures
  can reach strong sample efficiency. Applied to literature extraction,
  analogous object-centric, uncertainty-aware reasoning could enhance
  extraction quality by treating each paper's abstract as a structured
  observation to be parsed against hypothesis definitions with
  principled uncertainty quantification. The broader convergence between
  AIF and deep learning demonstrated by AXIOM
  \citep{heins2025axiom}---which outperforms DreamerV3 through
  principled active inference planning---further validates this
  trajectory.
\item
  \textbf{Domain adaptation.} The framework is domain-agnostic by
  design. Adaptation to foundation models, quantum computing, or
  synthetic biology requires only domain-specific hypothesis definitions
  and keyword lists within the A/B/C taxonomy.
\item
  \textbf{Energy-Based Model convergence.} Systematic cross-referencing
  of the Active Inference literature with the broader Energy-Based Model
  research program \citep{lecun2006tutorial}---including Helmholtz
  machines \citep{dayan1995helmholtz}, contrastive divergence training
  \citep{hinton2002training}, and variational autoencoders
  \citep{kingma2014auto}---would illuminate the mathematical convergence
  between variational free energy minimization in biological systems and
  energy-based learning in artificial systems. This analysis could
  reveal shared mathematical structures that are currently obscured by
  disciplinary siloing.
\item
  \textbf{Agent-Centric Architecture.} The deployment of structured
  \texttt{SKILL.md} files across the codebase establishes a Native Model
  Context Protocol (MCP) interface. This shifts the project from a
  static repository into an agentic workspace, empowering AI coding
  assistants to autonomously navigate, test, and expand the pipeline's
  capabilities while strictly adhering to localized infrastructural
  constraints.
\end{enumerate}
\end{block}
\end{frame}

\begin{frame}{Limitations}
\protect\phantomsection\label{limitations}
This meta-analysis is subject to three primary limitations that govern
the interpretation of its findings. First, as a deliberate design
trade-off for corpus-scale processing, we are only analyzing abstracts
for claims. This abstract-only extraction inherently misses granular
empirical evidence, methodological caveats, and nuanced findings
embedded within full texts. Second, our analysis possesses limited
bibliographic depth; keyword-based retrieval across disjoint APIs
produces a snapshot that inevitably contains coverage gaps, cross-source
identifier mismatches, and false negatives. Third, recognizing these
constraints, the current findings do not constitute a definitive, static
ledger of the field's evidence base. Rather, the work is presented as a
living open-source milestone and a computational scaffold for ongoing
collaboration. By publishing this architecture, we invite the community
to iteratively curate custom bibliographies, refine extraction prompts,
and gradually extend the system toward full-text analysis.
\end{frame}

\begin{frame}{Broader Impact}
\protect\phantomsection\label{broader-impact}
The vision motivating this work is straightforward: a living literature
review---a continuously updated knowledge graph tracking what a field
claims, what evidence supports those claims, and where the frontiers of
understanding lie. This vision builds on the foundation established by
Knight et al.~\citep{knight2022fep}, who identified the development of
systems that could ``encompass increased scope of relevant works,''
``integrate multiple forms of annotation and participation,'' and
``facilitate integration of manual and artificial contributions'' as key
goals for the field.

By demonstrating that LLM-driven assertion extraction can produce
scalable, queryable representations of scientific evidence---processing
\(N = 849\) papers spanning approximately two and a half decades
(2005--2026), extracting structured assertions, and evaluating 8 core
hypotheses---this work provides a reusable architecture for realizing
this vision. The corpus window begins in 2005 to capture Energy-Based
Model and variational Bayesian antecedents that predate the Free Energy
Principle label itself; the formal FEP was introduced in 2006
\citep{friston2006free} and reached its core elaboration by 2010
\citep{friston2010free}. The citation network metrics (1,678 edges,
0.23\% density, mean in-degree 2.0) characterize the field's structure,
which has grown at a 16.99\% CAGR while diversifying across 5
application domains.

The limitations of keyword-based retrieval across disjoint academic
repositories mean that any retrieved corpus will contain both false
positives and false negatives. There is no single threshold that
perfectly defines inclusion or exclusion for a dynamic,
interdisciplinary research field. The primary contribution of this work
is therefore not a definitive corpus but an open-source, modularly
updatable, and versioned software package. This tool is built in
reference to custom literature bibliographies that can be iteratively
curated for relevance by the community.

The combination of multi-source retrieval, LLM-based extraction, and
probabilistic knowledge graph construction provides a reusable template
that advances each of these goals. A complementary pathway is emerging
through Retrieval-Augmented Generation (RAG) architectures that ground
LLMs directly in knowledge graphs, reducing hallucination and enabling
real-time, context-aware reasoning over structured evidence
\citep{fan2024survey}. Integrating our nanopublication graph into such a
RAG system would enable natural-language querying of the evidence base,
further lowering the barrier for community engagement. The recent
release of nanopub-js v0.1.0 \citep{kuhn2026nanopubjs}---enabling
browser-based creation, signing, and querying of
nanopublications---lowers the barrier for community-contributed
assertions, bringing the participatory evidence curation envisioned by
Knight et al.~within practical reach. As LLM capabilities improve and
standardized metadata adoption grows, the cost of maintaining such
systems will decrease while their utility increases. By open-sourcing
the pipeline and publishing the schema, we provide both a concrete tool
for the Active Inference community and a modular blueprint that other
fields can adapt and refine.

\textbf{Data and code availability.} The pipeline source code,
configuration, and manuscript templates are available in the project
repository (see \texttt{metadata.repository} in \texttt{config.yaml} or
the manuscript front matter). Nanopublications are persisted as JSON
Lines (for incremental runs) and RDF/TriG (nanopub.net-compliant); both
can be archived with the code release or on a data repository (e.g.,
Zenodo) for citation and long-term access.

Community recommendations, actionable implications, and open questions
arising from this work are detailed in the
\hyperref[sec:discussion]{Discussion}.
\end{frame}

\end{document}
